\def\SECTION{2.1 Inductive Reasoning}
\documentclass[aspectratio=169,14pt]{beamer}
\usepackage[T1]{fontenc}
\usepackage{mathpazo}
\usepackage{amsmath,amssymb}
\usepackage{tikz}
\usetikzlibrary{arrows.meta,calc,decorations.pathreplacing}
\definecolor{fblue}{HTML}{1F4E79}
\definecolor{fteal}{HTML}{0F7C7C}
\definecolor{fgold}{HTML}{C9962B}
\definecolor{fred}{HTML}{B3261E}
\definecolor{fink}{HTML}{16233B}
\definecolor{fsoft}{HTML}{EEF3F8}
\definecolor{fgoldsoft}{HTML}{FBF1DC}
\setbeamertemplate{navigation symbols}{}
\setbeamercolor{normal text}{fg=fink}
\setbeamercolor{frametitle}{fg=fblue}
\setbeamerfont{frametitle}{series=\bfseries,size=\Large}
\setbeamercolor{structure}{fg=fteal}
\setbeamertemplate{itemize item}{\color{fteal}$\blacktriangleright$}
\setbeamertemplate{frametitle}{\vspace{6pt}\insertframetitle\par\vspace{-2pt}{\color{fgold}\rule{2.2cm}{2pt}}}
\setbeamertemplate{footline}{\hspace{0.6cm}{\color{fblue}\scriptsize\bfseries Fang\textcolor{fteal}{Math}}\hfill{\scriptsize\color{gray}\SECTION\quad\insertframenumber}\hspace{0.6cm}\vspace{6pt}}
\newcommand{\key}[1]{{\color{fteal}\textbf{#1}}}
\newcommand{\ans}[1]{\colorbox{fgoldsoft}{\color{fink}\ensuremath{#1}}}
\newcommand{\warn}[1]{\begin{beamercolorbox}[sep=8pt,rounded=true]{warnbox}{\color{fred}\textbf{Watch out!}}\enskip #1\end{beamercolorbox}}
\setbeamercolor{warnbox}{bg=fred!8,fg=fink}
\setbeamercolor{turnbox}{bg=fgoldsoft,fg=fink}
\setbeamercolor{defbox}{bg=fsoft,fg=fink}
\newcommand{\defn}[1]{\begin{beamercolorbox}[sep=8pt,rounded=true]{defbox}#1\end{beamercolorbox}}
\newenvironment{yourturn}{\begin{beamercolorbox}[sep=10pt,rounded=true]{turnbox}{\color{fgold}\Large\bfseries Your turn!}\quad{\small Pause the video and try it.}\par\medskip}{\end{beamercolorbox}}
\newcommand{\titleframe}[3]{%
\begin{frame}[plain]
\begin{tikzpicture}[remember picture,overlay]
\fill[fblue] (current page.north west) rectangle ([yshift=-0.9cm]current page.north east);
\node[anchor=west,text=white,font=\bfseries] at ([xshift=0.6cm,yshift=-0.45cm]current page.north west) {Fang\textcolor{fgold}{Math}};
\node[anchor=east,text=white,font=\small] at ([xshift=-0.6cm,yshift=-0.45cm]current page.north east) {Geometry $\cdot$ Chapter \CHNO};
\end{tikzpicture}
\vfill
{\color{fteal}\large\bfseries #1}\par\medskip
{\color{fblue}\fontsize{30}{34}\selectfont\bfseries #2}\par\medskip
{\color{fgold}\rule{3cm}{3pt}}\par\bigskip
{\large #3}
\vfill
{\small\color{gray} fangmath.com $\cdot$ YouTube @fangmathcom}
\end{frame}}
\newcommand{\closeframe}[1]{%
\begin{frame}{Remember}
#1
\vfill
\begin{center}{\color{fblue}\large\bfseries Now practice at your level:}\par\smallskip
{\Large\color{fteal}\bfseries fangmath.com} \enskip{\color{gray}$\to$ Geometry $\to$ Chapter \CHNO}\end{center}
\end{frame}}

\usepackage{pgfplots}\pgfplotsset{compat=1.17}
\newcommand{\grid}[4]{\draw[step=1,gray!25] (#1,#2) grid (#3,#4);\draw[-{Stealth},gray] (#1,0)--(#3+0.3,0) node[right]{\small$x$};\draw[-{Stealth},gray] (0,#2)--(0,#4+0.3) node[above]{\small$y$};}
\def\CHNO{2}
\usetikzlibrary{angles,quotes}
\begin{document}
\titleframe{Lesson 2.1}{Patterns \& Inductive Reasoning}{Look for the pattern, then predict}
\begin{frame}{What is inductive reasoning?}
\onslide<1->{\defn{\key{Inductive reasoning}: look at examples, find a \key{pattern}, and use it to \key{predict}.}}\par\medskip
\onslide<2->{\textbf{1.} Look \quad $\to$ \quad \textbf{2.} Describe the rule \quad $\to$ \quad \textbf{3.} Predict}\par\medskip
\end{frame}
\begin{frame}{A number pattern}
\onslide<1->{{\Large $3,\ 7,\ 11,\ 15,\ \ldots$}}\par\medskip
\onslide<2->{Rule: \key{add 4} each time \quad {\small\color{gray}(check every pair: $7-3$, $11-7$, $15-11$)}}\par\medskip
\onslide<3->{Next term: $15 + 4 = \ans{19}$}\par\medskip
\onslide<4->{10th term: $3 + 9 \cdot 4 = \ans{39}$ \quad {\small\color{gray}start at 3, add 4 nine times}}\par\medskip
\end{frame}
\begin{frame}{A picture pattern}
\onslide<1->{\tikz{\foreach \n/\x in {1/0,2/1.2,3/3,4/5.4}{\foreach \i in {1,...,\n}{\foreach \j in {1,...,\n}{\fill[fteal] (\x+\i*0.3,\j*0.3) circle(0.1);}}}}\par\smallskip $1,\ 4,\ 9,\ 16,\ \ldots$}\par\medskip
\onslide<2->{These are \key{square numbers}: $1^2, 2^2, 3^2, 4^2$. \quad Next: $5^2 = \ans{25}$}\par\medskip
\end{frame}
\begin{frame}{Conjecture and counterexample}
\onslide<1->{\key{Conjecture}: a statement you believe is true, based on a pattern.}\par\medskip
\onslide<2->{Conjecture: ``The sum of two numbers is always greater than each number.''}\par\medskip
\onslide<3->{\key{Counterexample}: $-2 + 5 = 3$, \ and $3 < 5$. \quad So the conjecture is \textbf{false}.}\par\medskip
\end{frame}
\begin{frame}{Watch out}
\onslide<1->{$2,\ 4,\ 8,\ \ldots$ \quad times 2? \ $\to 16$}\par\medskip
\onslide<2->{\warn{Or $+2, +4, +6$? \ $\to 14$. \ Two rules can fit the same first terms --- that's why inductive reasoning gives a \textbf{conjecture}, not a proof.}}\par\medskip
\end{frame}
\begin{frame}{}
\begin{yourturn}Find the next two terms: \quad {\Large $1,\ 3,\ 6,\ 10,\ \ldots$}\end{yourturn}\par\medskip
\onslide<2->{The differences are $+2, +3, +4$, so next $+5$, then $+6$: \quad $10 + 5 = \ans{15}$, \ $15 + 6 = \ans{21}$}\par\medskip
\end{frame}
\closeframe{\begin{itemize}\setlength\itemsep{8pt}
\item \key{Inductive reasoning}: look, describe the rule, predict.
\item Check the rule on \textbf{every} pair of terms.
\item A pattern gives a \key{conjecture}; one \key{counterexample} shows it is false.
\end{itemize}}
\end{document}
